\section{Introduction}
\label{intro}

The parton model invented by Feynman has been one of the most useful languages in describing the physics
of strong interactions at high energy~\cite{Feynman:1969ej}. Partons are idealized objects that exist when
hadrons travel at the speed of light. Thus in modern language, they are effective degrees of
freedoms in an effective theory, for instance, the soft-collinear
effective theory~\cite{Bauer:2000yr}.
The parton properties of a hadron, such as parton distributions, parton
wave functions or amplitudes, parton correlations, etc., are low-energy properties which
require solutions of quantum chromodynamics (QCD) in the strong coupling region.

When formulated in the rest frame of the hadron, parton physics represents
the space-time correlations of QCD quarks or gluons along the light-cone direction ($\xi^-=(t-z)/\sqrt{2}$),
which involves an explicit time dependence~\cite{Sterman:1994ce}.
(We use $\xi^\mu = (t, x, y, z)$ $(\mu=0,1,2,3)$ to denote space-time coordinates, although $x$, $y$, and $z$ are commonly used to
denote momentum fractions as well.)
In physics, the time-dependent
correlations are in a class by themselves because the time evolution calls for the presence of the
interaction-dependent Hamiltonian $H$, and the observables are called ``dynamical''.

Dynamical correlations are hard to calculate for at least two different reasons: 1)
if one inserts a set of intermediate Hamiltonian eigenstates to simplify the evolution,
one needs to know the wave functions of
all the intermediate excited states. Thus knowing just the initial (often ground) state wave function is not
sufficient to calculate the time-dependent correlations. Methods designed
to approximate the ground state might not be applicable to the excited states, and thus dynamical
correlations are usually calculated less reliably than time-independent or static
observables for which the initial state wave function is sufficient. 2) Since the time-dependence
involves the unitary evolution operator, $e^{-iHt}$, it is a problem for numerical Monte Carlo simulations as
it contains a sign-changing phase. In fact, the majority of the Monte Carlo approaches designed
to calculate static properties cannot be readily applied to dynamical correlations. There have been attempts
to calculate instead an imaginary or Euclidean evolution with $e^{-H\tau}$, and to analytically continue
the result to the real time~\cite{Nakahara:1999vy,Asakawa:2000tr}. However,
these approaches are hardly reliable because
the analytical functions are unknown. An alternative trick is to Taylor-expand the
time-evolution operator, obtaining the time dependence through the sum of an infinite series.
Such an approach again is very limited because of the poor precision of calculations
for higher moments~\cite{Hagler:2009ni}. For the above reasons, the time-dependent correlations
are generally referred to as ``Minkowskian'' in the sense that time and space correlations
are intrinsically different. On the other hand, the time-independent correlations may
be called ``Euclidean'', because they can be readily calculated with Monte Carlo simulations,
except that there may be additional sign problems.

For Minkowskian observables, the traditional approach is to use analytical approximations.
For instance, the so-called Schwinger-Dyson approach solves the time dependence
in terms of dynamical equations~\cite{Roberts:1994dr}. For parton physics, in the so-called light-front quantization approach,
the coordinates are redefined such that the new time is a combination of space and the old time, and
the light-cone correlations become time-independent~\cite{Burkardt:1995ct}. However, the resulting theory cannot be
solved readily using Monte Carlo methods because the metric is not positive definite.

Fortunately, there are certain ``dynamical'' problems that are not intrinsically ``Minkowskian'' in
the sense that with some clever formulation, they can be solved using Monte Carlo methods. For instance, scattering phase-shifts usually involve
complex S-matrix elements which cannot be calculated directly using a Monte Carlo approach.
However, one can put the theory in a box, calculating the bound state levels as a function of the boundary condition and extracting the phase shifts from bound state energy levels~\cite{Luscher:1986pf,Luscher:1990ck,Luscher:1990ux,Briceno:2017max}. A second example is certain Minkowskian correlators can be directly analytically continued to
Euclidean region without encountering any ``poles'', and therefore,
can be calculated directly in Euclidean space~\cite{Ji:2001wha}.

Recently it was suggested by the authors that parton physics in non-perturbative QCD belongs to a class
of problems that are not intrinsically Minkowskian~\cite{Ji:2013fga,Ji:2013dva,Hatta:2013gta,Xiong:2013bka,
Ji:2014gla,Ji:2014lra,Ji:2015jwa,Ji:2015qla,Xiong:2015nua}. This observation is based on two important
points: 1) light-cone correlators can be approximated by near-light-cone correlators through
QCD perturbation theory, thanks to asymptotic freedom. 2) the Lorentz boost is dynamical,
and it can be used to absorb the dynamical evolution between the fields through
boosting a hadron state. Therefore, all remainder dynamics is included in the boosted hadron wave function,
not in the probe operators. In this way, all parton physics can be accessed from equal-time correlations in a large-momentum hadron
state which are calculable in Euclidean Monte Carlo simulations.

The approach is an effective field theory in the following sense.
For any time-dependent light-cone operator $O(\xi^-_1,\xi^-_2, ...)$,
one can construct an infinite number of Euclidean quasi-operators
${\cal O}_i(z_1,z_2,..)$, where $i$ is a label for all such operators, and $z_i$ are the spatial coordinates
of the fields along the direction of hadron momentum, which is taken as the third coordinate direction.
The matrix element of such an operator in a hadron state with a large momentum $p_z$, calculated for example in lattice QCD,  has a large momentum expansion~\cite{Ji:2014gla}
\begin{equation} \label{eq:fact}
\langle p|{\cal O}_i(z_1,z_2,..)|p\rangle
= C_{i0}(\alpha_s,p_z)\otimes\langle p|O(\xi^-_1,\xi^-_2, ...)|p\rangle + {\cal O}(1/p^2_z)\ ,
\end{equation}
where we have omitted all the renomalization scales, and $C_{i0}$ is a hard coefficient. The higher-order
corrections can be calculated systematically.
Since in this effective theory the large momentum is used to approach the light-cone, it is called the large momentum effective field theory
or LaMET for short.

It is worthwhile to point out that the choice of the Euclidean quasi-operator above is not unique. Any quasi-operator that yields the same light-front physics in the infinite momentum limit can be used to probe the desired light-cone operator. All such quasi-operators form a universality class. In practice, the existence of the universality class allows one to make an optimized choice of the quasi-operators so that the ${\cal O}(1/p_z^2)$ correction can be made as small as possible. One such example is the calculation of the total gluon polarization
in the nucleon~\cite{Hatta:2013gta}.

Recently, there has been much interest about this theory
in the literature, with some questioning its validity or
effectiveness. Here we provide some discussions aiming at a further exposition of this approach.
In Ref.~\cite{Carlson:2017gpk}, the factorization within the LaMET approach has been questioned. Following Ref.~\cite{Briceno:2017cpo},
we explicitly show that the problem arises from the procedure of analytical continuation. In Ref.~\cite{Rossi:2017muf},
the issue of power divergences has been raised based on an analysis of moments. We show that this problem does not exist in the non-local approach
where the only power divergence arises from the self-energy of the Wilson lines.
Finally, in Refs.~\cite{Radyushkin:2017cyf,Orginos:2017kos}, a new approach has been proposed to extract
parton distributions.  We show that the Ioffe-time
distribution is an alternative way to extract
parton distribution from the same lattice observable, and hence is entirely equivalent to the LaMET approach.
Besides, the apparent scaling observed in that approach might or might not help with the convergence problem.
